%% 7. まとめ
\setband{\T{\S7 まとめ}{\S7 Summary}}

\begin{frame}{\T{決めてほしいこと（3 つ）}{Decisions needed (3)}}
\rl{\T{決定}{Decisions}}
\normalsize
\T{%
\begin{enumerate}
\item \textbf{論文の型}：「測定と説明」の型で書く（骨組みの項目 4 を前のページの案にする）か
\item \textbf{実験 4 を流すか}：見込みの得は +1 ポイント以下だが，「一様な重みが弱点か」を否定も含めて書ける（H200 1 台，4〜6 時間，約 50〜80 ドル）
\item \textbf{信頼度の筋を続けるか}：選択の結果は AFEW \N{sseldAFEWfive}〜\N{sseldAFEWtwofive}，MOSEI \N{sseldMOSEItwofive}〜\N{sseldMOSEIfive}．続けるなら，見込みの偏りを直す版（指数と見込みを別の抜き方で作る）を設計して検査に回す．止めるなら，信頼度の見積もりの結果だけを論文の探索の節に書く
\end{enumerate}
}{%
\begin{enumerate}
\item \textbf{Paper type}: write a ``measurement and explanation'' paper (skeleton item 4 as drafted on the previous page)?
\item \textbf{Run Experiment 4?} Expected gain ≤ +1 point, but it settles whether uniform weighting is a weakness, including a negative answer (one H200, 4--6 h, ≈\$50--80)
\item \textbf{Continue the trust-estimate line?} Selection result: AFEW \N{sseldAFEWfive} to \N{sseldAFEWtwofive}, MOSEI \N{sseldMOSEItwofive} to \N{sseldMOSEIfive}. If yes, design a version that removes the estimate bias (separate held-out draws for exponents and estimate) and send it through the checks. If no, report only the reliability-estimate result in the paper's exploration section
\end{enumerate}
}
\end{frame}
